% 00-liminaires.tex — Remerciements, résumé et abstract (hors pagination du corps).

\chapter*{Remerciements}
\addcontentsline{toc}{chapter}{Remerciements}

\begin{arediger}{½ page}
Promoteur, lecteurs, enseignants et étudiants du cours de Physique~1 ayant
participé aux tests, collègues du groupe.
\end{arediger}

\chapter*{Résumé}
\addcontentsline{toc}{chapter}{Résumé}

\begin{arediger}{1 page (résumé + abstract)}
Un paragraphe de 200 à 250 mots, écrit en dernier : problème (difficultés des
étudiants de BAC1 en mécanique : forces, mouvement circulaire, énergie),
démarche (site de simulations interactives développé avec une IA encadrée par
des règles, des tests et des agents ; mode jeu), principaux résultats
(chapitres~\ref{chap:resultats-quantitatifs} et~\ref{chap:resultats-qualitatifs}),
conclusion. Mots-clés : simulation interactive, mécanique, BAC1, ludification,
dynamique basée sur les positions, IA générative.
\end{arediger}

\section*{Abstract}

\begin{arediger}{traduction du résumé}
Version anglaise du résumé, mêmes mots-clés.
\end{arediger}
